% =====================================================================
%  Capítulo 11 — Entrenamiento
% =====================================================================
\chapter{Entrenamiento}

La preparación psicológica es la más importante: las cualidades físicas de nada sirven cuando no son controlables debido a las emociones intensas por la descarga de adrenalina (visión en túnel, aumento de latidos, rigidez muscular). Pero también es cierto que cuanto mejor sea el estado de nuestro cuerpo, más posibilidades habrá de salir bien librada.

Por ello se iniciará el entrenamiento básico de las armas corporales, descartando movimientos rebuscados y aprovechando la natural potencia de algunas zonas. Primaremos el uso de golpes, que pueden entrenarse en el aire (cosa no muy recomendable) pero que se dominan con un mínimo de equipo: manoplas, saco, o incluso un saco de dormir sujetado por alguien.

\begin{itemize}
  \item La efectividad no se consigue solamente con el conocimiento, sino con la práctica.
  \item La práctica ha de ser progresiva: primero lentamente, cuidando la forma correcta sin forzar el organismo; después incrementando velocidad y potencia.
  \item No golpeamos reteniendo la respiración ni inspirando, sino exhalando para acompañar el golpe.
  \item Antes de entrenar con fuerza, calienta el organismo unos quince minutos.
  \item Visualiza: imagina que estás frente a un adversario y en qué parte de su cuerpo golpeas.
  \item Si es posible, entrena con compañero y equipos protectores para acostumbrarte a la sensación de ser agredida, e incluye escenarios con voz: marcar límites gritando antes del contacto.
\end{itemize}
